\documentclass{article}
\usepackage{amsmath,amssymb,hyperref}
\newcommand{\Q}{\mathbb Q}\newcommand{\Z}{\mathbb Z}\newcommand{\N}{\mathbb N}
\begin{document}
\title{A bounded-step Bernstein bridge to rational first-survival slopes}
\date{16 November 2026: unregistered candidate corollary for audit}
\maketitle
This is a new written deduction for checking, conditional on the retained internal actual-module construction and formal Bernstein input. It is not independently verified or registered, and it does not solve exceptional Mathieu safety or pure-to-safe GVC. Its abstract lattice lemma is written separately from the differential-module application. No historical priority is claimed.

\section{The new leverage in the completed field proof}
For a characteristic-zero field $K$, a nonzero polynomial $P$ of degree $D>0$, and a nonzero formal symbol $\lambda\in(\xi)K[[\xi]]$, put $T=\lambda(D_x)$ and
\[
 \tau_k=\min\{j\ge0:T^kP^j\ne0\},\qquad
 v_k=\lambda(\xi)^k1.
\]
The retained proof \texttt{GENERIC\_BERNSTEIN\_EQUATIONS\_AND\_FORMAL\_GAP\_CLOSURE.md} constructs the actual holonomic image $M\subset K[x][[\eta]]$, with $\nu_\eta(v_k)=\tau_k$, and gives a single nonzero polynomial
\[
 p(k,\tau_k)=0\quad\hbox{for every sufficiently large $k$ with finite $\tau_k$}.
\]
No purity or safety is needed for that polynomial constraint. If some $\tau_k$ is infinite, every subsequent one is infinite. We consider the alternative where all are finite.

\section{An ordinary Bernstein equation bounds the increments}
Apply the primary generalized Bernstein theorem to the SAME actual holonomic $M$, its element $1$, and coefficient function $\lambda$. It gives
\[
 B(s)\lambda^s1=Q(s)\lambda^{s+1}1,
 \quad B\in K[s]\setminus0,\quad Q\in\mathcal D_R[s],
\]
in $M_\lambda[s]\lambda^s$, where $R=K[[\xi,\eta]]$. This is the ordinary one-parameter equation, not the second Bernstein equation used for the polynomial order constraint.

There is an actual unlocalized identity at every sufficiently large integer $k$:
\[
 B(k)v_k=Q(k)v_{k+1}.                                      \tag{1}
\]
Here is the localization check. Expand the finite operator by Leibniz and write the formal difference as $\lambda^{s-a}u(s)$ for one $a\ge0$ and $u(s)\in M[s]$. Its localized vanishing gives $\lambda^A u(s)=0$ for one uniform $A$, because there are finitely many coefficients. For $k\ge a+A$, the actual difference is $\lambda^{k-a}u(k)=0$. Thus (1) is an identity on the ORIGINAL powers, without replacing $k$ by $k+A$ or imposing eventuality on completed multipliers.

In PBW form let $J$ be the largest $\partial_\eta$ exponent in $Q(s)$. In the actual model, formal coefficients in $R$ preserve $\eta$-order, $\partial_{\xi_i}=-x_i$ preserves it, and $\partial_\eta=D_\eta+P$ lowers it by at most one. Excluding the finitely many integer roots of $B$, (1) yields
\[
 \tau_k\ge\tau_{k+1}-J.
\]
The original operator powers imply monotonicity. Hence eventually
\[
                  0\le\tau_{k+1}-\tau_k\le J.              \tag{2}
\]
In particular $\tau_k=O(k)$. This conclusion relies on the same actual-module and primary Bernstein hypotheses as the retained field proof; it does not revive the false scalar-time equation or product-Rees model.

Primary input: Mebkhout--Narvaez-Macarro, Ann. sci. ENS 24 (1991), Theorem3.1.1 and Section4.2, \href{https://doi.org/10.24033/asens.1627}{DOI 10.24033/asens.1627}. The cached primary PDF was already read in the targeted proof audit.

\section{Abstract lattice lemma: bounded steps force a rational affine cluster}
\textbf{Lemma.} Suppose integers $a_k$ satisfy, for all sufficiently large $k$,
\[
 0\le a_{k+1}-a_k\le J,\qquad p(k,a_k)=0,
 \quad 0\ne p\in\Q[X,Y].
\]
Then for some rational $\rho\ge0$,
\[
                         a_k=\rho k+O(1).                \tag{3}
\]
The sequence need not be eventually periodic after subtracting this slope.

\textbf{Proof.} Bounded steps give $a_k=O(k)$. Remove vertical and repeated polynomial factors and pass beyond the finite real discriminant values. The remaining real roots are finitely many algebraic graph branches $Y=f_i(X)$. Their convergent Puiseux expansions at infinity show that each branch is eventually either $O(X)$ or outside the range of $a_k$. Discard the latter branches. The surviving branches have bounded derivatives. For each pair, the same expansions show that $f_i-f_j$ is bounded or tends in absolute value to infinity.

Group the branches by bounded difference. Distinct groups cannot exchange the selected branch at consecutive late indices: their separation tends to infinity, while the derivative bounds and (2) bound every one-step change. Thus one fixed group supplies $a_k$ for EVERY sufficiently large integer $k$.

We show that this group contains a rational affine line. A nonlinear algebraic graph branch has eventually nonzero second derivative of constant sign; otherwise its algebraic second derivative would have infinitely many zeros and vanish identically, making the branch affine. It is therefore strictly convex or concave on a final interval. Its arc length over $[k_0,N]$ is $O(N)$ by the derivative bound.

Here is the needed elementary lattice estimate. Order its $L_0$ integer points along such a strictly convex/concave arc. Consecutive chord vectors are distinct nonzero elements of $\Z^2$, since their slopes are strictly monotone. There are $O(R^2)$ integer vectors of norm at most $R$. At least a fixed fraction of $L_0-1$ distinct vectors therefore have norm at least $c\sqrt{L_0-1}$. Their summed lengths are at least $c'(L_0-1)^{3/2}$ and at most the arc length. Consequently $L_0=O(1+N^{2/3})$. This is the elementary Jarnik bound; the argument is given here for arcs, so no closed-curve hypothesis is being silently imported.

Each nonrational affine line contains at most one integer point: two such points determine a rational slope and intercept. If the fixed group had no rational affine line, its finitely many nonlinear branches and nonrational lines would supply only $O(N^{2/3})$ points with $k\le N$. They must supply $N-O(1)$ points, a contradiction. Choose its rational line $\rho X+c$. Every other branch in the same group differs from it by $O(1)$, proving (3). Nonnegativity of $a_k$ gives $\rho\ge0$. \hfill$\square$

For provenance of the lattice bridge see Jarnik, \emph{Uber die Gitterpunkte auf konvexen Kurven}, Math. Z.24 (1926),500--518, \href{https://www.dml.cz/handle/10338.dmlcz/500690}{primary archive}. Convergent Puiseux expansions of plane algebraic curves at infinity are the classical Newton--Puiseux input. No Bombieri--Pila bound or integer-point search is required.

If the original polynomial constraint has coefficients in arbitrary $K$, expand its finitely many coefficients in a basis of their $\Q$-span. At integer pairs every rational component vanishes. A nonzero rational component supplies the polynomial used by the lemma, exactly as in the retained field proof.

\section{Candidate all-symbol consequences over fields}
Combining (1)--(3) with the retained polynomial constraint gives, in the all-finite alternative,
\[
               \boxed{\tau_k=\rho k+O(1),\quad\rho\in\Q_{>0}.} \tag{4}
\]
Positivity follows because $T^k$ has order at least $rk$, with $r=\operatorname{ord}\lambda>0$, and $P^j$ has degree $Dj$; hence $\tau_k\ge rk/D$.

Define the ACTUAL heights $h_m=\max\{k:T^kP^m\ne0\}$ and $H_m=\max_{0\le j\le m}h_j$. We do NOT assume $T^kP^j$ stays nonzero for every $j\ge\tau_k$. Instead its definition gives exactly
\[
 \tau_k\le m\ \Longleftrightarrow\ \exists j\le m:\ h_j\ge k,
 \qquad H_m=\max\{k:\tau_k\le m\}.
\]
Therefore (4) implies
\[
 H_m=\frac m\rho+O(1),\quad
 \limsup_m\frac{h_m}m=\frac1\rho,\quad
 \liminf_m\frac{m-h_m}m=1-\frac1\rho.                    \tag{5}
\]
The limsup equality follows since $H_m\ge h_m$, while any eventual bound $h_j\le(L+\epsilon)j$ gives $H_m\le\max(C,(L+\epsilon)m)$.

$H_m$ is also the exact height for the passive suspension $P+u$, since its different powers of the ignored variable $u$ prevent cancellation. Thus the rational bounded-error law applies to the suspended heights. Only the limsup statement, not a bounded-error formula, is claimed for the original unsuspended $h_m$.

Under ALL-ORDER purity, $\rho\ge1$. The retained formal commutator/gap equivalence then gives: in the all-finite case, safety is equivalent to $\rho>1$. At $\rho=1$, $\tau_k-k$ is bounded, which contradicts the necessary divergent shifted deficit for safety. The nilpotent alternative $\tau_N=\infty$ is safe by the same commutator estimate. This is a classification in terms of an actual asymptotic invariant, not a proof every pure input has $\rho>1$.

Given the two Bernstein certificates and their orders, a further constructive project is to isolate which rational line-cluster is selected and extract certified numerical bounds. This note implements no general producer for Bernstein certificates of arbitrary infinite formal input, computes no least multiplier cutoff, and claims no universal numerical $C$. No all-ring rational-slope theorem follows here.

\section{Controls and limits}
For $P=xy$ and $T=\partial_x^N+\partial_y^M$, $N,M\ge2$, distinct terms cannot cancel and
\[
 \tau_k=\min_{0\le i\le k}\max(Ni,M(k-i)),\qquad
 \rho=\frac{NM}{N+M}.
\]
Choosing an integer nearest $Mk/(N+M)$ gives bounded error. The jump is nonnegative and at most $\max(N,M)$ (add either derivative summand to an optimal term). For $M=N+1$, the reduced slope denominator is $2N+1$, unbounded even at fixed phase degree two. Thus a bound on possible denominators depending only on phase degree is false. $N\ge2,M=N+1$ gives $\rho>1$ and all-order purity by $1/N+1/M<1$.

Two abstract controls show what the new step adds. The sequence $a_k=k+1$ has slope1 and is compatible with a genuine unsafe pure seed; a rational slope is not a safety theorem. The sequence $a_k=k+\mathbf1_{\{k\text{ a square}\}}$ is monotone with jumps at most2 and satisfies $(a_k-k)(a_k-k-1)=0$, but its bounded intercept is not eventually periodic. Thus (3) cannot assert eventual periodicity.

The earlier nonmonotone square-index polynomial-curve control has unbounded jumps and is excluded by (2); the false product-law and raw scalar-time routes remain retired. The full exceptional Mathieu target is not replaced by this corollary.
\end{document}
